% Universidade Federal de Pernambuco
% Departamento de Engenharia Elétrica
% Autora: Andrea Maria Cavalcanti   (andrea.marianogueira@ufpe.br)
% Autor: Victor Rodrigo Leite       (victor.rodrigo@ufpe.br)  
% Recife, 2022

% Utilize o comando \section para criar uma seção
% Utilize o comando \subsection para criar subseção

% Capítulo Introdução

\chapter{Introdução}
\label{chap:introducao}

\textcolor{red}{[Contextualização do trabalho, incluindo sua relevância/justificativa]}








\section{Objetivos}

\subsection{Geral}

\subsection{Específicos}

\section{TESTE GLOSSÁRIO}
	Este é um teste de paragrafo, escreverei bastante para a quebra de linha automática do glossarioe etou difitando um novo paragrado teste de quebra de linha, abredito que já esta tudo certo por Latex é muito bom


\section{Metodologia}


\section{Organização Textual}

Este trabalho de conclusão de curso é organizado da seguinte forma:

\begin{itemize}
    \item \textbf{Capítulo 2:} 
    
    \item \textbf{Capítulo 3:}

    \item \textbf{Capítulo 4:};

    \item \textbf{Capítulo 5:} 
\end{itemize}



